\label{chap:brewrc}

Commercial ReRAM memories implement undocumented error-correcting codes
that impose deterministic, polarity-dependent timing perturbations on
every write operation. These perturbations dominate the observable write
latency distributions of deployed devices, yet no non-invasive technique
existed to recover the underlying ECC structure prior to this work. This
chapter presents BREW-RC (Bit-exact Recovery of ECC from Write Timing in
ReRAM Crossbars), a technique that recovers the bit-exact ECC parity
submatrix of a commercial ReRAM memory using only write-access timing,
without physical probing or device modification. The chapter makes three
contributions. First, it identifies and experimentally characterizes an
undocumented timing side-channel in commercial ReRAM devices arising
from ECC parity-bit switching activity. Second, it introduces BREW-RC,
which formulates ECC recovery as a SAT problem whose constraints are
derived from the polarity structure of observed write latency
distributions. Third, it demonstrates that the recovered parity
submatrix accurately predicts write-latency distributions in commercial
devices, enabling precise timing models for side-channel analysis and
reliability characterization. Section~\ref{sec:brewrc-setup} describes
the experimental platform and measurement methodology.
Section~\ref{sec:brewrc-investigation} presents a series of
architectural write latency experiments that characterize the timing
side-channel and motivate the BREW-RC approach.
Section~\ref{sec:brewrc-method} presents the BREW-RC technique in
detail. Section~\ref{sec:brewrc-results} presents timing model
construction and validation results.
